%----------
%	DATA BALANCING
%----------	
As highlighted by \cite{liebowitz-2024}, hate speech detection faces significant challenges due to "limited available datasets, which are often highly imbalanced and exhibit topical and lexical biases". This challenge was also present in our study, as discussed in Section~\ref{target_variables}, where we observed significant class imbalances across the three different types of classification tasks we aimed to tackle. We could have ignored this fact and proceeded with the experimentation but that wouldn't be the ideal thing to do as it would highly penalize the minority classes in terms of the different metrics considered, in particular F1 Score, precision and recall.

Moreover, \cite{rathpisey-2019} provide other different data resampling techniques, such as \textit{Random Over-sampling} (ROS),  \textit{Synthetic Minority Oversampling Technique} 
(SMOTE) and \textit{Adaptive Synthetic} (ADASYN). According to their findings, all oversampling techniques led to improved results across various models, whereas downsampling was effective only for the naive Bayes model. We did not explore downsampling further because this technique reduces the number of instances in the majority class to match the minority class, which would have resulted in a significantly reduced dataset size. Given that our dataset is already limited, applying downsampling would have left us with too few examples.

\begin{itemize}
    \item \textbf{Over-sampling:} In this case, the foundation behind this idea is to populate the minority class with examples so it reaches the number of examples of the majority class. The different approaches that were considered are the same as presented in \cite{rathpisey-2019}. This technique was applied only to the training set after the split to avoid duplication of test examples, including also shuffling to prevent order bias.
        \begin{itemize}
            \item \textbf{Random Oversampling (ROS):} The way to proceed in this case is to generate copied examples of the minority class and randomly select them with replacement until the classes are balanced. In our experiments, we refer to this technique as \textit{"upsampling"}.
            
            \item \textbf{Synthetic Minority Oversampling Technique (SMOTE):} This technique was first introduced in \cite{chawla-2002}. The idea is to artificially generate new examples (synthetic examples) in the minority class, rather than simply oversampling with replacement. This method operates by first identifying the k-nearest neighbours for each sample in the minority class within the feature space, reflecting the similarities among instances. A random neighbour is then chosen, and a new synthetic sample is generated through linear interpolation between the chosen neighbour and the original sample. Specifically, the new sample $x_{new}$ is calculated using the formula:
            \begin{equation}
                x_{\text{new}} = x_{\text{original}} + \lambda \cdot (x_{\text{neighbor}} - x_{\text{original}})
            \end{equation}
            
            The python library \textit{imblearn} provides a built-in function for this task.
            
            \item \textbf{Adaptive Synthetic (ADASYN):} This idea was first presented in \cite{he-2008}. It extends the methodology of SMOTE by introducing a weighted distribution mechanism to handle imbalanced datasets more effectively. Unlike SMOTE, which generates the same number of synthetic samples for each instance in the minority class, ADASYN adjusts the number of synthetic samples based on the local density of the class. It assigns weights to the minority class samples depending on how difficult they are to classify; those that are harder to learn (i.e., more likely to be misclassified) receive higher weights. 
            
            The python library \textit{imblearn} also provides a built-in function for this task.
        \end{itemize}
\end{itemize}


